\documentclass[11pt]{article}
\usepackage[letterpaper,margin=1in]{geometry}
\usepackage[T1]{fontenc}
\usepackage[utf8]{inputenc}
\usepackage{lmodern,microtype}
\usepackage{amsmath,amssymb,mathtools,bm,slashed}
\usepackage{xcolor,enumitem,booktabs,array,needspace}
\usepackage[most]{tcolorbox}
\usepackage{fancyhdr,titlesec,parskip,hyperref}
\definecolor{qftblue}{HTML}{245A8D}
\definecolor{qftgreen}{HTML}{2F6B4F}
\definecolor{qftgray}{HTML}{F4F6F8}
\hypersetup{colorlinks=true,linkcolor=qftblue,urlcolor=qftblue,pdfauthor={Nima Lashkari}}
\pagestyle{fancy}\fancyhf{}
\lhead{PHYS 662: Quantum Field Theory I}\cfoot{\thepage}
\setlength{\headheight}{14pt}
\titleformat{\section}{\Large\bfseries\color{qftblue}}{\thesection.}{0.6em}{}
\titleformat{\subsection}{\large\bfseries\color{qftgreen}}{\thesubsection}{0.6em}{}
\setlist[enumerate]{leftmargin=2em,itemsep=0.25em,topsep=0.35em}
\setlist[itemize]{leftmargin=2em,itemsep=0.2em,topsep=0.35em}
\newcommand{\dd}{\mathrm d}
\newcommand{\ii}{\mathrm i}
\newcommand{\eps}{\varepsilon}
\newcommand{\tr}{\operatorname{tr}}
\newcommand{\Tr}{\operatorname{Tr}}
\newcommand{\diag}{\operatorname{diag}}
\newcommand{\Lag}{\mathcal L}
\newcommand{\D}{\mathcal D}
\newcommand{\vac}{\lvert0\rangle}
\newcommand{\vbra}{\langle0\rvert}
\newcommand{\Tord}{\mathrm T}
\newcommand{\ket}[1]{\lvert#1\rangle}
\newcommand{\bra}[1]{\langle#1\rvert}
\newtcolorbox{keyidea}[1][]{enhanced,breakable,colback=qftgray,colframe=qftblue,
 boxrule=0.7pt,arc=1.5mm,title={#1},fonttitle=\bfseries,coltitle=white}
\newtcolorbox{remarkbox}[1][]{enhanced,breakable,colback=white,colframe=qftgreen,
 boxrule=0.6pt,arc=1.5mm,title={#1},fonttitle=\bfseries,coltitle=white}
\newcommand{\module}[1]{\clearpage\section{#1}}
\newcommand{\lecturefront}[4]{%
\rhead{Lecture #1}\hypersetup{pdftitle={PHYS 662 Lecture #1: #2}}
\begin{document}\pagenumbering{gobble}
\begin{titlepage}\centering\vspace*{1.05in}
{\scshape\Large PHYS 662 --- Quantum Field Theory I\par}
\vspace{0.35in}{\Huge\bfseries\color{qftblue}Lecture #1\par}
\vspace{0.2in}{\LARGE\bfseries #2\par}
\vspace{0.2in}{\large\color{qftgreen}#3\par}
\vspace{0.6in}{\large Nima Lashkari\par}{\large Fall 2026\par}
\vfill\begin{minipage}{0.86\textwidth}\small #4\par\medskip
Course page: \url{https://nimalashkari.com/quantum-field-theory-2026-part-i/}
\end{minipage}\vspace{0.5in}\end{titlepage}
\pagenumbering{roman}\setcounter{tocdepth}{1}\tableofcontents
\vfill\begin{remarkbox}[Place in the course]
Lectures 17--18 develop the quantum-mechanical construction. Lecture 19
extends Gaussian integration and source methods to scalar fields.
Lecture 20 introduces Grassmann integration and fermionic fields.
Problem Set 6 brings these four steps together.
\end{remarkbox}\clearpage\pagenumbering{arabic}}
\lecturefront{18}{Path Integrals in Quantum Mechanics II}
{Perturbations, scattering, and Euclidean time}
{The free kernel becomes a building block for interactions. We then rotate
to imaginary time, where evolution projects onto the ground state and a
trace becomes a thermal partition function. These constructions supply
the boundary conditions needed for vacuum field theory.}

\section{Perturbing the free kernel}
Keep \(\hbar\) explicit and write \(H=H_0+V\), with
\(H_0=\bm p^2/(2m)\). Lecture 17 fixed the normalization of
\(K_0(\bm q_f,\bm q_i;T)\). Expand the interaction part of the weight:
\begin{equation}
 e^{-\frac{\ii}{\hbar}\int_{t_i}^{t_f}V(q(t))\dd t}
 =1-\frac{\ii}{\hbar}\int\dd t\,V(q(t))
 +\frac1{2!}\left(-\frac{\ii}{\hbar}\right)^2
 \int\dd t\,\dd t'\,V(q(t))V(q(t'))+\cdots.
\end{equation}
Composition cuts the free path at each insertion. In one dimension,
\begin{align}
 K^{(1)}(q_f,t_f;q_i,t_i)
 =-\frac{\ii}{\hbar}\int_{t_i}^{t_f}\dd t\int\dd q\,
 K_0(q_f,t_f;q,t)V(q)K_0(q,t;q_i,t_i).                \label{eq:bornkernel}
\end{align}
The insertion transfers amplitude between free propagation segments.
For the second-order term, restrict the integration to \(t_1<t_2\):
\begin{align}
 K^{(2)}={}&\left(-\frac{\ii}{\hbar}\right)^2
 \int_{t_i}^{t_f}\dd t_2\int_{t_i}^{t_2}\dd t_1
 \int\dd q_2\dd q_1\,
 K_0(q_f,t_f;q_2,t_2)V(q_2)\nonumber\\[-2pt]
 &\hspace{17mm}\times K_0(q_2,t_2;q_1,t_1)V(q_1)
 K_0(q_1,t_1;q_i,t_i).
\end{align}
The two possible time orderings cancel the \(2!\). At order \(n\),
either integrate over ordered times with no factorial, or over the
whole time cube with time ordering and a factor \(1/n!\).

Equivalently, the interaction-picture operator obeys
\begin{equation}
 U_I=\Tord\exp\!\left[-\frac{\ii}{\hbar}\int\dd t\,V_I(t)\right],
 \qquad V_I(t)=e^{\ii H_0t/\hbar}Ve^{-\ii H_0t/\hbar}.
\end{equation}
This is the Dyson expansion derived by slicing histories. It is not
necessary that \(\int V\dd t\ll\hbar\) on every path. Perturbation
theory is an expansion of amplitudes; its accuracy depends on couplings,
energies, duration, and possible resonances. Long times can generate
large terms even when the potential is weak.

\begin{remarkbox}[Two bodies and a fixed target]
For a potential \(V(\bm r_1-\bm r_2)\), center-of-mass and relative
motion separate. The relative kernel uses the reduced mass
\(\mu=m_1m_2/(m_1+m_2)\). A fixed heavy target is the limit
\(m_2\gg m_1\); momentum need not be conserved for the probe alone.
\end{remarkbox}

\module{Energy space and the resolvent}
Let \(T=t_f-t_i\). The unit-source retarded inverse from Lecture 17 is
\(G_R(T)=\theta(T)e^{-\ii HT/\hbar}/(\ii\hbar)\).
Use the energy Fourier transform
\begin{equation}
 \widetilde G_R(E)=\int_{-\infty}^{\infty}\dd T\,
 e^{\ii ET/\hbar}G_R(T),\qquad
 G_R(T)=\int\frac{\dd E}{2\pi\hbar}
 e^{-\ii ET/\hbar}\widetilde G_R(E).
\end{equation}
For an energy eigenvalue \(E_n\), damping the positive-time integral
by \(e^{-\varepsilon T/\hbar}\) gives
\begin{equation}
 \frac1{\ii\hbar}\int_0^\infty\dd T\,
 e^{\ii(E-E_n+\ii\varepsilon)T/\hbar}
 =\frac1{E-E_n+\ii\varepsilon}.
\end{equation}
Thus
\begin{equation}
 R^+(E)\equiv\widetilde G_R(E)=(E-H+\ii0)^{-1}.        \label{eq:resolvent}
\end{equation}
The Fourier transform of \(\theta(T)K\), in contrast, is \(\ii\hbar\)
times this resolvent. The prescription is fixed by positive-time
evolution; writing a formal inverse without it leaves boundary
conditions unspecified.

For a free particle,
\begin{equation}
 R_0^+(E;\bm q_f,\bm q_i)=\int\frac{\dd^d p}{(2\pi\hbar)^d}
 \frac{e^{\ii\bm p\cdot(\bm q_f-\bm q_i)/\hbar}}
 {E-\bm p^2/(2m)+\ii0}.
\end{equation}
The distribution identity
\(1/(x+\ii0)=\operatorname{PV}(1/x)-\ii\pi\delta(x)\)
separates the off-shell principal value from the on-shell part.
The intermediate momentum is integrated freely; it is not constrained
to satisfy \(E=\bm p^2/(2m)\) at every insertion.

We can now organize the perturbation theory in $V$ for the resolvent using the operator identity $(R^+)^{-1}-(R^+_0)^{-1}=(R^+_0)^{-1}(R^+_0-R^+)(R^+)^{-1}$. Iterating this by substituting $R^+$ into the right-hand-side gives
\begin{align}
 R^+&=R_0^++R_0^+VR^+\\
 &=R_0^++R_0^+VR_0^++R_0^+VR_0^+VR_0^++\cdots.
                                                        \label{eq:resborn}
\end{align}
This is the energy-space version of the ordered kernel expansion.
Each integration over an intermediate time has become another
resolvent, and the interaction insertions remain in their operator order.
Note that in the time domain, the total elapsed time $T$ is shared among all the propagation segments: $T=(t_f-t_n)+\cdots+(t_2-t_1)+(t_1-t_i)$ which is a convolution in time. By the convolution theorem, the convolution in time becomes multiplication in Fourier (energy) space.

\begin{keyidea}[Boundary conditions are part of an inverse]
\((E-H)R=1\) has several distributional solutions on the continuous
spectrum. The \(+\ii0\) in \eqref{eq:resolvent} selects the one obtained
from retarded evolution, and becomes the outgoing-wave condition in
stationary scattering.
\end{keyidea}

\module{From evolution to a scattering amplitude}
For a localized, sufficiently short-range potential, scattering compares
interacting motion with freely moving wave packets far before and after
the encounter. The asymptotic free phases must be removed:
\begin{equation}
 S=\lim_{\substack{t_f\to+\infty\\t_i\to-\infty}}
 e^{\ii H_0t_f/\hbar}e^{-\ii H(t_f-t_i)/\hbar}
 e^{-\ii H_0t_i/\hbar}.                              \label{eq:smatrix}
\end{equation}
This limit is understood on the scattering subspace, with wave packets
or an adiabatic switching prescription. Long-range Coulomb scattering
needs a modification of the simplest free asymptotic description.
A bare infinite-time limit of the Schr\"odinger-picture kernel is not
the scattering operator.

Use \(\langle\bm p|\bm p'\rangle=(2\pi\hbar)^d
\delta^{(d)}(\bm p-\bm p')\). At first order,
\begin{align}
 S_{fi}^{(1)}&=-\frac{\ii}{\hbar}
 \int_{-\infty}^{\infty}\dd t\,
 e^{\ii(E_f-E_i)t/\hbar}\langle\bm p_f|V|\bm p_i\rangle\\
 &=-2\pi\ii\,\delta(E_f-E_i)\widetilde V(\bm p_f-\bm p_i),\\
 \widetilde V(\bm Q)&=\int\dd^d q\,
 e^{-\ii\bm Q\cdot\bm q/\hbar}V(\bm q).
\end{align}
The time integral is \(2\pi\hbar\delta(E_f-E_i)\), so the
\(\hbar\) cancels the vertex factor. A static target conserves the
probe energy while absorbing momentum. The angular dependence of the
first Born amplitude samples spatial Fourier components of the potential.

Introduce a transition operator \(\mathcal T(E)\), distinct from the
time-ordering symbol \(\Tord\):
\begin{align}
 \mathcal T(E)&=V+VR_0^+(E)V+\cdots=V+VR_0^+(E)\mathcal T(E),\\
 S_{fi}&=(2\pi\hbar)^d\delta^{(d)}(\bm p_f-\bm p_i)
 -2\pi\ii\delta(E_f-E_i)
 \langle\bm p_f|\mathcal T(E_i)|\bm p_i\rangle.       \label{eq:transition}
\end{align}
For example, its second term contains
\begin{equation}
 \int\frac{\dd^d k}{(2\pi\hbar)^d}
 \frac{\widetilde V(\bm p_f-\bm k)\widetilde V(\bm k-\bm p_i)}
 {E_i-\bm k^2/(2m)+\ii0}.
\end{equation}
This denominator encodes propagation between insertions. It does not
describe an observable intermediate particle violating energy
conservation. The measurable statement is the full on-shell amplitude.

\module{Imaginary time and the heat kernel}
Continue a positive real-time interval to \(T=-\ii\tau\), with
\(\tau>0\). Spectrally, the oscillatory factor becomes a damping factor:
\begin{equation}
 K_E(q_f,q_i;\tau)=\bra{q_f}e^{-H\tau/\hbar}\ket{q_i}
 =\sum_n\psi_n(q_f)\psi_n(q_i)^*e^{-E_n\tau/\hbar}.  \label{eq:spectral}
\end{equation}
For a continuous spectrum the sum includes an integral. Assuming the
Hamiltonian is bounded below, this evolution is a well-defined
semigroup. It is not unitary and need not preserve the norm of a state.

With \(t=-\ii\tau\), \(\dot q=\ii q'\), so
\begin{equation}
 \frac{\ii S}{\hbar}=-\frac{S_E}{\hbar},\qquad
 S_E[q]=\int\dd\tau\left[\frac m2(q')^2+V(q)\right].
                                                        \label{eq:euclaction}
\end{equation}
The potential changes sign relative to the real-time Lagrangian;
it is added to the Euclidean kinetic term. The regulated kernel is
\begin{equation}
 K_E=\lim_{N\to\infty}\left(\frac{m}{2\pi\hbar a}\right)^{N/2}
 \int\prod_{j=1}^{N-1}\dd q_j\,
 e^{-\frac1\hbar\sum_j[\frac{m(q_{j+1}-q_j)^2}{2a}+aV(q_j)]},
 \qquad a=\frac{\tau}{N}.
\end{equation}
For a real scalar potential the weight is nonnegative, when the
regulated integral exists. The action need not itself be positive:
shifting \(V\) by a constant changes the energy zero and multiplies
the kernel by an overall exponential.

Differentiating the operator definition gives a heat-type equation,
\begin{equation}
 \hbar\partial_\tau K_E
 =\left(\frac{\hbar^2}{2m}\partial_{q_f}^2-V(q_f)\right)K_E,
 \qquad K_E(0)=\delta(q_f-q_i).
\end{equation}
In particular,
\begin{equation}
 K_{E,0}(q_f,q_i;\tau)
 =\left(\frac{m}{2\pi\hbar\tau}\right)^{1/2}
 e^{-m(q_f-q_i)^2/(2\hbar\tau)}.
\end{equation}
Its spreading is diffusive, with variance \(\hbar\tau/m\). Imaginary
time is an analytic and computational tool; it is not a second
physical time coordinate. In interacting continuum theories the
analytic continuation and the regulator require care, rather than a
purely formal substitution into divergent expressions.

\module{Thermal traces and periodic bosonic paths}
Let \(\beta=1/(k_BT_{\rm phys})\). The thermal partition function is
\begin{equation}
 Z_\beta=\Tr e^{-\beta H}
 =\int\dd q\,K_E(q,q;\beta\hbar)
 =\int_{q(\beta\hbar)=q(0)}\D q\,e^{-S_E[q]/\hbar}.  \label{eq:thermal}
\end{equation}
The starting point of the loop is integrated. Fixing it would give a
diagonal matrix element, not a trace. The circumference of the
Euclidean time circle is \(\beta\hbar\). Periodic Fourier modes have
\(\omega_n=2\pi n/(\beta\hbar)\).

\subsection{An oscillator checks both the measure and the trace}
For \(V(q)=m\omega^2q^2/2\), the classical equation of motion is
\begin{align}
q''-\omega^2 q=0
\end{align}
whose solutions are $\sinh$ and $\cosh$. The Euclidean classical path with fixed
endpoints is
\begin{equation}
 q_{\rm cl}(s)=
 \frac{q_i\sinh[\omega(\tau-s)]+q_f\sinh(\omega s)}{\sinh(\omega\tau)}.
\end{equation}
Using its equation of motion, the on-shell action is the boundary
term \(\frac m2[qq']_0^\tau\). The fluctuation determinant is
endpoint independent. Fixing it by the short-time free limit gives
\begin{align}
 K_{E,\omega}(q_f,q_i;\tau)
 &=\sqrt{\frac{m\omega}{2\pi\hbar\sinh(\omega\tau)}}\,
 e^{-S_{E,\rm cl}/\hbar},\\
 S_{E,\rm cl}&=\frac{m\omega}{2\sinh(\omega\tau)}
 [(q_f^2+q_i^2)\cosh(\omega\tau)-2q_fq_i].
\end{align}
To determine the prefactor, insert \(K_E=\mathcal N(\tau)e^{-S_{E,\rm cl}/\hbar}\)
into the heat equation. The classical Hamilton--Jacobi equation cancels
the terms independent of \(\hbar\), leaving
\[
 \frac{\mathcal N'}{\mathcal N}
 =-\frac1{2m}\partial_{q_f}^2S_{E,\rm cl}
 =-\frac\omega2\coth(\omega\tau).
\]
Thus \(\mathcal N\propto[\sinh(\omega\tau)]^{-1/2}\); the short-time
free limit fixes its constant to the value displayed above.

Setting \(q_f=q_i=q\) and doing one ordinary Gaussian integral yields
\begin{equation}
 Z_\beta=\frac{1}{2\sinh(\beta\hbar\omega/2)}
 =\frac{e^{-\beta\hbar\omega/2}}{1-e^{-\beta\hbar\omega}}.
                                                        \label{eq:hoz}
\end{equation}
This agrees with \(\sum_{n\geq0}e^{-\beta\hbar\omega(n+1/2)}\),
including the ground-state energy. A normal-ordered Hamiltonian
\(\hbar\omega a^\dagger a\) would instead omit the numerator.

\begin{keyidea}[What the classical statistical analogy means]
A regulated Euclidean path integral resembles a statistical sum with
local couplings between neighboring time slices. In field theory this
adds one Euclidean coordinate to the spatial lattice. At finite
temperature that coordinate is a circle with circumference $\beta$. This analogy does not identify real-time dynamics
with a classical equilibrium process.
\end{keyidea}

\module{Vacuum projection and Feynman boundary conditions}
Suppose there is a unique ground state \(\ket0\) and a gap above it.
For a trial state with \(c_0=\langle0|\chi\rangle\neq0\),
\begin{equation}
 e^{-H\tau/\hbar}\ket\chi
 =e^{-E_0\tau/\hbar}
 \left[c_0\ket0+\sum_{n>0}c_ne^{-(E_n-E_0)\tau/\hbar}\ket n\right].
\end{equation}
Taking the limit $\tau\to \infty$ and dividing by the appropriate normalization projects onto the vacuum.
If the vacuum is degenerate, it projects onto the ground-state
subspace, and the surviving state depends on the boundary data.

For \(\tau_2>\tau_1\), a normalized insertion formula is
\begin{equation}
 \lim_{L\to\infty}
 \frac{\bra\chi e^{-H(L-\tau_2)/\hbar}q
 e^{-H(\tau_2-\tau_1)/\hbar}q
 e^{-H(\tau_1+L)/\hbar}\ket\chi}
 {\bra\chi e^{-2LH/\hbar}\ket\chi}
 =\langle0|q_E(\tau_2)q_E(\tau_1)|0\rangle,
                                                        \label{eq:projection}
\end{equation}
where \(q_E(\tau)=e^{H\tau/\hbar}qe^{-H\tau/\hbar}\).
The denominator cancels the vacuum energy and the trial-state overlap.
It also explains why vacuum generating functionals are conventionally
normalized to equal one at zero source.

Real-time projection can be implemented by taking a long interval
along a slightly downward-tilted contour:
\begin{equation}
 e^{-\ii H T(1-\ii\epsilon)/\hbar}
 =e^{-\ii HT/\hbar}e^{-\epsilon HT/\hbar},\qquad\epsilon>0.
\end{equation}
Subtract \(E_0\) if convenient. First send \(T\to\infty\) at fixed
\(\epsilon\), then remove the tilt. States with energy above \(E_0\)
are suppressed. The resulting vacuum insertions are time ordered.

For the harmonic oscillator,
\begin{equation}
 C_F(t)=\frac{\hbar}{2m\omega}e^{-\ii\omega|t|}
 =\int\frac{\dd\nu}{2\pi}\,
 \frac{\ii\hbar\,e^{-\ii\nu t}}{m(\nu^2-\omega^2+\ii0)}.
\end{equation}
For \(t>0\), the positive-frequency pole lies below the real axis;
for \(t<0\), the negative-frequency pole lies above it. These are
Feynman boundary conditions. They differ from the retarded inverse
of a first-order Schr\"odinger operator in Section 2.

The endpoint kernel, the vacuum generating functional, and the thermal
trace use closely related actions but different boundary data. Keeping
these three constructions distinct prevents many normalization errors.

\module{Euclidean correlations and the energy gap}
Insert an energy basis into a vacuum two-point function of a Hermitian
operator \(O\). For \(\tau>0\), its connected part is
\begin{equation}
 C_{E,c}(\tau)=
 \sum_{n>0}|\langle0|O|n\rangle|^2
 e^{-(E_n-E_0)\tau/\hbar}.                           \label{eq:gap}
\end{equation}
The large-\(\tau\) decay measures the smallest excitation energy to
which \(O\) couples. A symmetry can make the overlap with the first
excited state vanish. It is therefore more precise to speak of the
gap in this operator's channel. A continuum of states can add a power
of \(\tau\) multiplying the exponential.

For the oscillator coordinate, only the first excited state contributes:
\begin{equation}
 C_E(\tau)=\frac{\hbar}{2m\omega}e^{-\omega|\tau|},\qquad
 m(-\partial_\tau^2+\omega^2)C_E(\tau)=\hbar\delta(\tau).
\end{equation}
The derivative jump at zero fixes the normalization. Fourier
transformation gives \(\widetilde C_E(\nu)=\hbar/[m(\nu^2+\omega^2)]\).
This is a positive inverse of the Euclidean quadratic operator.

\subsection{Gaussian sources anticipate field theory}
Write \(A_E=m(-\partial_\tau^2+\omega^2)\) and use vacuum boundary
conditions at infinite Euclidean time. Define
\begin{equation}
 Z_E[J]=\frac{\int\D q\,e^{-\frac1{2\hbar}qA_Eq+
                                    \frac1\hbar Jq}}
 {\int\D q\,e^{-\frac1{2\hbar}qA_Eq}}
 =\exp\!\left[\frac1{2\hbar}JA_E^{-1}J\right].
                                                        \label{eq:gaussiansource}
\end{equation}
Here products include time integrations. Completing the square by
\(q=q'+A_E^{-1}J\) proves the second equality. Derivatives with
\(\hbar\delta/\delta J\) insert \(q\), so
\(C_E=\hbar A_E^{-1}\). In particular,
\begin{equation}
 \langle\Tord_E q_1q_2q_3q_4\rangle
 =C_{12}C_{34}+C_{13}C_{24}+C_{14}C_{23}.
\end{equation}
All higher Gaussian moments reduce to pairings. The logarithm of
\(Z_E\) is quadratic, so connected moments of order greater than two
vanish, even though the full even moments do not.

At finite temperature the same differential operator acts on periodic
functions on a circle. Its inverse is a periodic image sum of
\(C_E\), not the infinite-line kernel above. Boundary conditions
therefore remain part of the definition even for a Gaussian integral.

\module{Euclidean saddles and tunneling}
Euclidean variation gives
\begin{equation}
 m q''=V'(q),\qquad
 \frac m2(q')^2-V(q)=\text{constant}.                \label{eq:euclmotion}
\end{equation}
This is formally motion in the inverted potential. It is not the
real-time Newton equation. For two degenerate minima \(q_\pm\) with
\(V(q_\pm)=0\), a finite-action interpolating saddle on the infinite
line has \(q'\to0\) at both ends. Its first integral is
\begin{equation}
 q'=\sqrt{\frac{2V(q)}{m}},\qquad
 S_I=\int_{q_-}^{q_+}\dd q\,\sqrt{2mV(q)}.           \label{eq:instanton}
\end{equation}
The reversed solution has the opposite sign for \(q'\) and the same
action. Such Euclidean saddles are called instantons and anti-instantons.
They provide a controlled starting point for weak tunneling when
\(S_I/\hbar\gg1\).

For the quartic double well \(V(q)=\lambda(q^2-a^2)^2\) which has degenerate vacua at $q=\pm a$ and the potential barrier $V(0)=\lambda a^4$, ordinary
perturbation theory about one minimum describes local oscillations.
Classically, a particle at zero energy sits at either vacuum, and never learns that the other well exists. However, quantum mechanically, there is an amplitude for tunneling. When the particle in one well has energy $E<V(q)$, it still cannot go over the barrier, however, tunneling describes going through the potential barrier, using the language of instantons.
Paths connecting the wells add contributions proportional to
\(e^{-S_I/\hbar}\). The prefactor requires a fluctuation determinant.
Translations of the instanton center are a zero mode; its integration
must replace a naive zero eigenvalue in that determinant.

Over a long Euclidean time, the particle does not tunnel just once, it can tunnel back and forth many times. In a dilute-instanton regime, since the probability for tunneling is exponentially suppressed, the events are rare, and independent (well-separated compared to the width of a single instanton). In this limit, the leading splitting has the form
\begin{equation}
 \Delta E\simeq2\hbar\kappa e^{-S_I/\hbar},
\end{equation}
where \(\kappa\) has units of inverse time and contains the fluctuation
factor. Summing many well-separated transitions is needed to obtain
the long-time spectrum. Neither the instanton exponential alone nor
one finite-time saddle is the exact double-well path integral. Problem
Set 6 asks for the trajectory and action, and supplies this framework
without requiring the determinant calculation.

\subsection*{Questions to check understanding}
\begin{enumerate}
\item Why does restricting insertion times remove a factorial?
\item Which factors of \(\hbar\) enter the energy delta function in
\eqref{eq:transition}?
\item Why must free evolution be removed before taking a scattering limit?
\item What changes between fixed endpoints, a vacuum projection, and a trace?
\item How do the short-time and large-time limits test the oscillator kernel?
\item Why can a Euclidean two-point function miss the lowest excitation?
\end{enumerate}
The scalar field will be a collection of coupled oscillators, with
the same Gaussian, boundary-condition, and source machinery.
\end{document}
